% Standalone insertion prepared by proof_audit/verify_cyclo_radix6.py.
% This file is not input by the main manuscript.
\begin{lemma}[Rigorous quadratic-splitting approximate-strong sampler]
\label{lem:radix-six-approx-strong}
Let
\[
  R=\mathbb Z[X]/(X^{128}+1),\qquad
  q=447183309836853377,
\]
and let $\mathcal D_6\subset R_q$ be the distribution obtained by sampling
the $128$ coefficients independently and uniformly from
$\mathcal A=\{-3,-2,-1,0,1,2\}$.  Then
$|\mathcal D_6|=6^{128}$ and
\[
  \Pr_{s,s'\leftarrow\mathcal D_6}
  \bigl[s-s'\notin R_q^\times\mid s\ne s'\bigr]
  \le
  \kappa_{\rm nu}^{(6)}
  :=\frac{64\cdot6^{-44}-6^{-128}}{1-6^{-128}}.
\]
Moreover, every $s\in\operatorname{supp}(\mathcal D_6)$ satisfies
$\|s\|_{\rm op}\le384$.
\end{lemma}

\begin{proof}
The accompanying machine-checkable Pocklington certificate proves that $q$
is prime.  It also verifies
\[
 q\equiv129\pmod {256},\qquad
 6^{64}\equiv-1\pmod q,\qquad
 \operatorname{ord}_q(6)=128.
\]
Consequently, writing $\omega_u=6^u$ for odd $u\in\{1,\ldots,127\}$,
\[
 X^{128}+1=\prod_{\substack{1\le u<128\\u\text{ odd}}}(X^2-\omega_u)\pmod q.
\]
There are $64$ distinct factors.  Each is irreducible: $v_2(q-1)=7$ and an
element of order $128$ is a nonsquare in $\mathbb F_q$.

Fix one factor $X^2-\omega_u$ and use the basis $\{1,Y\}$ of its quadratic
residue field.  The residue of
$s=\sum_{i=0}^{127}s_iX^i$ is $A_u+YB_u$, where
\[
 A_u=\sum_{j=0}^{63}s_{2j}\omega_u^j,
 \qquad
 B_u=\sum_{j=0}^{63}s_{2j+1}\omega_u^j.
\]
For $t\in\{0,\ldots,21\}$, let $r_t$ represent $u^{-1}t$ modulo $128$ and
put
\[
 (j_t,\epsilon_t)=
 \begin{cases}
 (r_t,+1),&r_t<64,\\
 (r_t-64,-1),&r_t\ge64.
 \end{cases}
\]
Then $\omega_u^{j_t}=\epsilon_t6^t$, and the $22$ indices $j_t$ are
distinct.  After fixing all other even coefficients, the unfixed part of
$A_u$ is therefore a translate of
$\sum_{t=0}^{21}\epsilon_td_t6^t$ with independent uniform
$d_t\in\mathcal A$.

This $22$-digit map is injective.  Indeed, the integer difference of two
colliding blocks has absolute value at most
\[
 5\sum_{t=0}^{21}6^t=6^{22}-1<q.
\]
Modular equality is thus integer equality.  Reduction modulo $6$, followed by
division by $6$ and induction, forces all digit differences in $[-5,5]$ to
vanish.  Hence $\max_a\Pr[A_u=a]\le6^{-22}$.  The identical argument for
$B_u$ uses disjoint odd coefficients, so coefficient independence gives
\[
 \max_{a,b}\Pr[(A_u,B_u)=(a,b)]\le6^{-44}.
\]
For independent $s,s'$, collision in this slot has probability at most
$6^{-44}$.  A difference is a nonunit only if such a collision occurs in one
of the $64$ CRT slots, whence
$\Pr[s-s'\notin R_q^\times]\le64\cdot6^{-44}$.
Since $\Pr[s=s']=6^{-128}$ and equality is contained in the nonunit event,
subtracting this repeat event and conditioning on $s\ne s'$ gives the stated
$\kappa_{\rm nu}^{(6)}$.

Finally, negacyclic convolution gives
$\|ts\bmod(X^{128}+1,q)\|_\infty
 \le\|t\|_\infty\|s\|_1$, while
$\|s\|_1\le128\cdot3=384$.
\end{proof}

\begin{corollary}[Explicit Cyclo interface]
\label{cor:radix-six-cyclo-interface}
In the challenge-dependent terms of Cyclo Theorem~3, the distribution
$\mathcal D_6$ may be instantiated with
\[
 |\mathcal D_6|=6^{128},\qquad
 \gamma=384,\qquad
 \kappa_{\rm nu}=\kappa_{\rm nu}^{(6)}.
\]
Thus the repeat and nonunit contributions are, respectively,
\[
 L6^{-128}
 \quad\text{and}\quad
 L\frac{64\cdot6^{-44}-6^{-128}}{1-6^{-128}}.
\]
Numerically,
$-\log_2\kappa_{\rm nu}^{(6)}\ge107.738350031731$.
\end{corollary}

\begin{remark}[Scope and trade-off]
Lemma~\ref{lem:radix-six-approx-strong} gives an explicit heuristic-free
quadratic-splitting replacement for the reasoning of Cyclo Lemma~9.  It does
not preserve Cyclo's original benchmark: the displayed modulus has $59$ bits
and is outside the original $50$-bit AVX-512/IFMA parameter line.  It does not
give $128$-bit nonunit security and does not improve the separate exact-strong
flagship radius line.  For this radix argument, reaching a $128$-bit
$64$-slot bound first occurs at $26$ digits, which requires
$q>6^{26}-1$ and hence at least a $68$-bit modulus; this is a construction-
specific trade-off, not a general lower bound.
\end{remark}
